\documentclass{CjC}
\usepackage[utf8]{inputenc}
\usepackage[T1]{fontenc}
\usepackage{geometry}
\usepackage{graphicx}
\usepackage{float}
\usepackage{hyperref}
\usepackage[english,indonesian]{babel}

\title{Optimasi Pemrosesan Template LaTeX Menggunakan Pendekatan Berbasis AI}
\author{Academic Support Agent}
\date{\today}

\begin{document}
\maketitle

\begin{abstract}

\end{abstract}

\chapter{PENDAHULUAN}


\chapter{TINJAUAN PUSTAKA}


% Original content (commented out):
% %% This is file `CjC_template_tex.tex',
%% is modified by Lingjia Dai based on the template 
%% provided by Chinese Journal of Computers (http://cjc.ict.ac.cn/).
%%
%% This version is capable with Overleaf (XeLaTeX).
%%
%% Update date: 2023/11/20

%% Import packages
% \documentclass[10.5pt,compsoc,UTF8]{CjC}
\documentclass[15.5pt]{article}
\usepackage{geometry}
\geometry{a4paper,left=2.54cm,right=2.54cm,top=3.17cm,bottom=3.17cm}
\usepackage{ctex}
\usepackage{xeCJK}
\usepackage{fontspec} 
% \setCJKmainfont{SimSun} %或\setCJKmainfont{KaiTi}
% \setCJKmonofont{SimSun} % 中文 SimSun
% \setmainfont{Times New Roman} % 英文 Times New Roman
\usepackage{graphicx}
\usepackage{subfigure}
\usepackage{multirow}
\usepackage{multicol}
\usepackage{amsmath,amsthm,amsfonts,amssymb}
\usepackage{color}
\usepackage{natbib}
\usepackage{makecell}
\usepackage{hyperref}
\usepackage{cleveref}
\crefname{figure}{图}{}
\crefname{table}{表}{表}
\crefname{algorithm}{算法}{算法}


\linespread{2.0} % 2倍行距
% \renewcommand{\CJKglue}{\hskip 20pt}
\ziju{0.11}

%===============================%
% format of footnote
\usepackage{scrextend}
\deffootnote{0em}{0em}{\thefootnotemark\quad}
\makeatletter
\def\blfootnote{\gdef\@thefnmark{}\@footnotetext}
\makeatother


\makeatletter
\newcommand{\rmnum}[1]{\romannumeral #1}
\newcommand{\Rmnum}[1]{\expandafter\@slowromancap\romannumeral #1@}
\makeatother

% format of different section level
\usepackage{titlesec}
\titleformat{\section}{\normalfont\normalsize\bfseries}{\thesection\,}{1em}{}
\titleformat{\subsection}{\normalfont\normalsize\bfseries}{\thesubsection\,}{1em}{}
\titleformat{\subsubsection}[runin]{}{\thesubsubsection}{1em}{}[]

% format of figure and table
\usepackage[font=small]{caption} % footnotesize
\usepackage{bicaption}
% \captionsetup[figure][bi-second]{name=Fig.}
\captionsetup{labelformat=default,labelsep=space} % or period

% table \hline new define
\makeatletter
\def\hlinewd#1{%
\noalign{\ifnum0=`}\fi\hrule \@height #1 %
\futurelet\reserved@a\@xhline}
\makeatother

\newcommand{\tr}[1]{\textcolor{red}{#1}}



\begin{document}
% \hyphenpenalty=50000 % 断字比例

%% Title
{\centering
{\fontsize{22pt}{22pt}\selectfont 
中文标题 （SimSun 22pt）
}

\vskip 5mm
作者$^{1}$，作者$^{2,3}$，作者$^{3}$

\vspace {5mm}
｜（1. 一级单位 ~ 二级单位，省份 ~ 城市 ~ 邮编；2.……）｜
\par % end centering
}

\vskip 5mm

\noindent 
{\bf 摘~~要：}
采用报道性摘要形式，内容包括：\textcircled{1} 目的(What you want to do(直接给出研究目的))；\textcircled{2} 方法(How you did it(详细陈述过程和方法))；\textcircled{3} 结果和结论(What results did you get and what conclusion can you draw(全面罗列结果和结论))；\textcircled{4}文章的创新之处(What is original in your paper)。摘要一般250字左右，采用第三人称的写法，尽量不使用“本人”~“本文”~“作者”~“笔者”等作为主语，不分段，不标注文献号，尽量不用公式和物理量字母等。

\noindent
\textbf{关键词：}列出3～5个关键词，请用“；”分开。根据学科、标题和摘要标引关键词，不用缩略词，写出中文全称。

\noindent
\textbf{中图分类号：}（投稿前请检索，交叉学科可列多个）\quad \textbf{文献标志码：A}
\vspace {5mm}


{\centering
{\fontsize{22pt}{22pt}\selectfont 
Title
}

\vskip 5mm
\textit{Author}$^{1}$，\textit{Author}$^{2,3}$，\textit{Author}$^{3}$

\vspace {5mm}
｜（1. Institute Name,  University Name, City Postcode, Province, China；2.……）｜

\par % end centering
}

\noindent
\textbf{Abstract：}
与中文摘要内容对应。用语简练，少用多个of句型；推荐使用现在时态叙述文章内容（方法、过程及结果结论）。

\noindent
\textbf{Key words：}
与中文关键词一一对应。如有英文缩略关键词，请先写全称，括号中加缩略词。

\blfootnote{
收稿日期：2023-00-00 修回日期：2023-00-00 录用日期：2023-00-00 \\
基金项目：国家自然科学基金（基金编号），上海市自然科学基金（基金编号）资助项目 \\
作者简介：姓名（1990—），博士生，从事……研究。 \\
通信作者：姓名，教授，博士生导师，电话（Tel.）：021-00000000；E-mail：XXXX。
}


\newpage

\begin{table}[]
  \caption{}
  \resizebox{1\columnwidth}{!}{
  \begin{tabular}{ccccccccccccccc}
  \hline
  \multicolumn{2}{c}{\multirow{3}{*}{\textbf{结合面类型}}} & 
  \multirow{3}{*}{\textbf{\makecell[c]{结合面最\\小宽度\\L/mm}}} & 
  \multicolumn{12}{c}{\textbf{最大间隙/mm}} \\
  & &
  & \multicolumn{3}{c}{$V<=100(\text{cm}^3)$} 
  & \multicolumn{3}{c}{$100<V<=500(\text{cm}^3)$} 
  & \multicolumn{3}{c}{$500<V<=2000(\text{cm}^3)$} 
  & \multicolumn{3}{c}{$V>2000(\text{cm}^3)$} \\
  & & & \Rmnum{1} & \Rmnum{2}A & \Rmnum{2}B & \Rmnum{1} & \Rmnum{2}A & \Rmnum{2}B & \Rmnum{1} & \Rmnum{2}A & \Rmnum{2}B & \Rmnum{1} & \Rmnum{2}A & \Rmnum{2}B \\
  \hline
  \multicolumn{2}{c}{\multirow{4}{*}{\textbf{\makecell[c]{平面结合\\面、圆筒形\\结合面或止\\口结合面}}}} 
  &   6    & 0.3  & 0.3 & 0.2 & ...  & ... & ... & ...  & ...  & ...  & ... & ... & ... \\
  & & 9.5  & 0.35 & 0.3 & 0.2 & 0.35 & 0.3 & 0.2 & 0.08 & 0.08 & 0.08 & ... & ... & ... \\
  & & 12.5 & 0.4  & 0.3 & 0.2 & 0.4  & 0.3 & 0.2 & 0.4  & 0.3  & 0.2  & 0.4 & 0.2 & 0.15 \\
  & & 25   & 0.5  & 0.4 & 0.2 & 0.5  & 0.4 & 0.2 & 0.5  & 0.4  & 0.2  & 0.2 & 0.4 & 0.2 \\
  \multirow{10}{*}{\textbf{\makecell[c]{旋转\\电机\\转轴\\结合\\面}}} 
  & \multirow{5}{*}{\textbf{\makecell[c]{滑动\\轴承}}} 
  &   6    & 0.3  & 0.3  & 0.2  & ...  & ... & ...  & ... & ... & ...  & ... & ... & ... \\
  & & 9.5  & 0.35 & 0.3  & 0.2  & 0.35 & 0.3 & 0.2  & ... & ... & ...  & ... & ... & ... \\
  & & 12.5 & 0.4  & 0.35 & 0.25 & 0.4  & 0.3 & 0.2  & 0.4 & 0.3 & 0.2  & 0.4 & 0.2 & ... \\
  & & 25   & 0.5  & 0.4  & 0.3  & 0.5  & 0.4 & 0.25 & 0.5 & 0.4 & 0.25 & 0.5 & 0.4 & 0.2 \\
  & & 40   & 0.6  & 0.5  & 0.4  & 0.6  & 0.5 & 0.3  & 0.6 & 0.5 & 0.3  & 0.6 & 0.5 & 0.25 \\
  %
  & \multirow{5}{*}{\textbf{\makecell[c]{滚动\\轴承}}} 
  &   6    & 0.45 & 0.45 & 0.3  & ...  & ... & ... & ... & ... & ... & ... & ... & ... \\
  & & 9.5  & 0.5  & 0.45 & 0.35 & 0.5  & 0.4 & 0.25 & ... & ... & ... & ... & ... & ... \\
  & & 12.5 & 0.6  & 0.5  & 0.4  & 0.6  & 0.45 & 0.3 & 0.6 & 0.45 & 0.3 & 0.6 & 0.3 & 0.2 \\
  & & 25   & 0.75 & 0.6  & 0.45 & 0.75 & 0.6 & 0.4 & 0.75 & 0.6 & 0.4 & 0.75 & 0.6 & 0.3 \\
  & & 40   & 0.8  & 0.75 & 0.6  & 0.8  & 0.75 & 0.45 & 0.8 & 0.8 & 0.45 & 0.8 & 0.75 & 0.4 \\
  \hline
  \end{tabular}
  }
  \label{tab:my-table}
\end{table}

\begin{center}
  \begin{table}[ht]
  \caption{Effects of the threshold in LocNet-HS.}
  \label{tab:ablation_threshold}
  \resizebox{1\columnwidth}{!}{
  \begin{tabular}{ccccccccc}
  \hline
  \multirow{2}{*}{\textbf{\makecell[c]{Number of\\Point Sources}}} & \multicolumn{2}{c}{\textbf{93\%}} 
  & \multicolumn{2}{c}{\textbf{95\%}}   
  & \multicolumn{2}{c}{\textbf{97\%}}
  & \multicolumn{2}{c}{\textbf{100\%}} \\
  & \textbf{Precision}  & \textbf{Recall}  
  & \textbf{Precision}  & \textbf{Recall} 
  & \textbf{Precision}  & \textbf{Recall} 
  & \textbf{Precision}  &\textbf{Recall}\\ 
  \hline
  5 & 96.85\% & 99.60\% & 99.33\% & \textbf{100.00\%} & \textbf{99.50\%} & \textbf{100.00\%} & 98.80\% & 99.60\% \\
  10 & 95.53\% & 98.90\% & 98.19\% & \textbf{99.10\%} & \textbf{98.34\%} & \textbf{99.10\%} & 97.22\% & 98.20\% \\
  15 & 96.46\% & 99.07\% & \textbf{96.74\%} & \textbf{99.20\%} & 96.58\% & 99.07\% & 96.23\% & 98.67\% \\
  20 & 94.82\% & 98.20\% & \textbf{95.17\%} & 98.10\% & 94.96\% & \textbf{98.50\%} & 95.07\% & 97.80\% \\
  25 & 93.37\% & \textbf{97.72\%} & \textbf{94.48\%} & 97.60\% & 93.80\% & 97.68\% & 93.87\% & 97.00\% \\
  30 & 92.64\% & 97.10\% & \textbf{94.37\%} & 97.47\% & 93.74\% & \textbf{97.50\%} & 94.01\% & 96.70\% \\
  35 & 90.87\% & 95.91\% & \textbf{92.15\%} & 96.14\% & 91.76\% & \textbf{96.57\%} & 92.00\% & 95.66\% \\
  40 & 90.31\% & \textbf{95.25\%} & \textbf{90.89\%} & 95.20\% & 90.43\% & 95.13\% & 90.21\% & 94.98\% \\
  45 & 89.11\% & 94.20\% & 89.14\% & \textbf{94.51\%} & 88.22\% & 94.33\% & \textbf{89.15\%} & 93.89\% \\
  \hline
  Average & 93.33\% & 97.33\% & \textbf{94.49\%} & 97.48\% & 94.15\% & \textbf{97.54\%} & 94.06\% & 96.94\%\\
  \hline
  \end{tabular}
  }
  \end{table}
  \end{center}

直接开始引言部分内容，不用标题。应开门见山，言简意赅。内容包括研究的目的、意义、主要方法、范围和背景(引文献)等，不要与摘要雷同或成为摘要的注释，避免公式推导和一般性方法介绍。最后一定要给出主要创新点。引言中一般不列图、表与公式。

\section{一级标题}

\subsection{二级标题}

\subsubsection{三级标题}
接正文内容。正文文字采用五号。


\subsection{二级标题}
插图(\Cref{fig:ex})应有自明性，切忌与表及文字表达重复。图题应简洁明确，但不宜选用“实验设备图”“函数关系图”等过于泛指的图名，须中英文对照，英文首字母大写，其余小写（专用词除外）。图中内容的注释用中文直接注于图中。线条图清晰规范，带标值的坐标轴须有完整的标目（量与单位名称）、标值线和对应标值、坐标端点标值；照片可以用彩图，比例说明采用标尺（Scale bar）。

\begin{figure}[ht]
  \centering
  \includegraphics[width=0.9\linewidth]{figs/example.jpg}
  % Version 1
  \caption{中文标题}
  \addtocounter{figure}{-1}
  \vspace{-13pt}
  \renewcommand{\figurename}{Fig.}
  \caption{Symbol Definitions}
  % % Version 2
  % \bicaption{中文标题}{English title}
  \label{fig:ex}
\end{figure}
\vspace{-15pt}


\subsection{二级标题}
表格(\Cref{tab:ex})应精心设计，具有自明性，切忌与图及文字表达重复。一律使用三线表，表格中同一列数值具有相同的量与单位，通栏排列，表题中英文对照，英文首字母大写，其余小写(专用词除外)。

\begin{table}[!h]
  \label{tab:ex}
  \centering
  \caption{中文标题}
  \addtocounter{table}{-1}
  \vspace{-13pt}
  \renewcommand{\tablename}{Tab.}
  \caption{Symbol definitions}
  \begin{tabular}{ccccc}
  % \hline
  \hlinewd{1.2pt}
  \multirow{2}{*}{\makecell[c]{网格类型}}
  & \multirow{2}{*}{\makecell[c]{最小横向\\尺寸/mm}}
  & \multirow{2}{*}{\makecell[c]{最小纵向\\尺寸/mm}}
  & \multirow{2}{*}{\makecell[c]{网格总量}}
  & \multirow{2}{*}{\makecell[c]{计算\\时间/h}}  
  \\
  & & & & \\
  \hlinewd{1pt}
  % 网格类型 & 最小横向尺寸/mm & 最小纵向尺寸/mm & 网格总量 & 计算时间/h \\
  % \hline
  Mesh1 & 0.8 & 0.48 & $1.7\times10^4$ & 1.0 \\
  Mesh2 & 0.8 & 0.48 & $1.7\times10^4$ & 1.0 \\
  Mesh3 & 0.8 & 0.48 & $1.7\times10^4$ & 1.0 \\
  \hlinewd{1.2pt}
  \end{tabular}
\end{table}



\subsection{二级标题}
并非每个公式均需编号，但后文中要引用的公式和重要的结论性公式必须编号。全文连续编号。公式的一般推导过程可不编号。

所有量符号均用单个字母表示（25个特征数符号除外），不宜用多个字母表示一个量，必要时可用角标区分；首次出现的量符号均要解释说明；不要重复使用同一字符代表不同变量；一般变量符号用斜体，如 $a$；矩阵、向量、张量等多维量符号用黑斜体， 如 $\boldsymbol{b}$, $\boldsymbol{C}$；一般下标表示缩写则正体；圆周率$\pi$、自然底数$\mathrm{e}$、微分符号$\mathrm{d}$、虚部$\mathrm{i}$，请正体。

\begin{equation}
  \frac{\partial u_i}{\partial t} + \frac{\partial}{\partial x_j}(u_i U_j) = 
  -\rho \frac{\partial p}{\partial x_i} + v \frac{\partial}{\partial x_j}
  (
    \frac{\partial u_i}{\partial x_j} + \frac{\partial u_i}{\partial x_i}
  )
  + \frac{\partial}{\partial x_j} (-\bar{u_i^{'}u_j^{'}})
\end{equation}

式中：$u_i$和$u_j$为速度时均量； $x_i$和$x_j$为坐标；$t$为时间；$\rho$为流体密度；$p$为压力时均量；$\frac{\partial}{\partial x_j} (-\bar{u_i^{'}u_j^{'}})$为雷诺应力。

\section{一级标题}

\section{一级标题}

\section{结论（结语）}
\noindent
稿件审理结果将在4个月内通知作者，个别稿件可能送审时间较长。逾期未收到稿件审理结果的请及时向编辑部查询。稿件录用后通知作者交审稿费版面费等。

（1）

（2）

（3）

$i j$

$\mathbb{P}$

文献著录格式（\textbf{\tr{作者姓氏全部大写，名字缩写或首字母大写双名中间不加连字符；中文文献需写出相应的英文}}）欢迎引用学报近年来发表的成果\cite{dAutonomousSearch2020}。


% \section*{参考文献}
% \bibliographystyle{gbt7714-numerical}
\bibliographystyle{plain}
\bibliography{ref.bib}
\end{document}




\end{document}